\documentclass[11pt]{article}
\usepackage[utf8]{inputenc}
\usepackage{amsmath,amssymb,amsthm}
\usepackage[margin=1in]{geometry}
\usepackage{hyperref}
\usepackage{xcolor}
\usepackage{enumitem}
\usepackage{newunicodechar}
\newunicodechar{ℝ}{\ensuremath{\mathbb{R}}}
\newunicodechar{ℕ}{\ensuremath{\mathbb{N}}}

\title{Tide retrospective:\\\emph{2D numerator asymptotic for the quartic-sextic potential (Z1)}}
\author{Daniel Murfet}
\date{18 May 2026}

\begin{document}
\maketitle

\begin{abstract}
The cross-project tide-candidate survey labelled the V2 follow-up
``numerator-side asymptotic''. This Tide discharges it: the missing
literal interpretation of the original survey wording, about the
unnormalised 2D integral $\int\!\int x^{2j} y^{2k}\,e^{-tL}\,dx\,dy$
scaling as $t^{-(2j+1)/4 - (2k+1)/6}$. The packaging mirrors V2:
an exact reformulation, a rescaled limit, and an asymptotic-
equivalence form.\footnote{The latter two via \texttt{Tendsto}
and \texttt{Asymptotics.IsEquivalent} respectively.} The proof is
one factorisation through separability plus two 1D moment closed
forms plus the same \texttt{rpow} algebra V2 used. 157 lines, single-attempt clean build,
with zero \texttt{sorry}, axiom, or \texttt{native\_decide}. Mirror-
shaped tide; GPT consult skipped per the tide skill's preflight, as
the candidate is structurally identical to V2 and the closed form was
verified numerically at ten test points to relative error
${\sim}10^{-16}$.
\end{abstract}

\section{Setting}

Z1 is the literal completion of the V2 story: V2 packaged T1's exact
moment closed form into asymptotic form; Z1 does the same for the
\emph{unnormalised numerator integral}. The two objects differ by the
2D partition function $\int\!\int e^{-tL}\,dx\,dy$, which itself
scales as $t^{-1/4 - 1/6} = t^{-5/12}$; the moment scaling
$t^{-(j/2 + k/3)}$ and the numerator scaling $t^{-((2j+1)/4 +
(2k+1)/6)}$ differ by exactly $t^{-5/12}$, as expected.

Two retrospective-time observations made Z1 cheaper than the survey
estimated:

\begin{itemize}[leftmargin=1.5em]
\item The 2D-numerator factor lemma already exists in the seabed.
  The survey direction guessed ``likely needs to be proved
  alongside, by direct factorisation of the 2D integral''. In fact,
  for any \texttt{addSeparable} potential $U + V$ and any separable
  observable, the 2D integral factors as
  $\int\!\int f(x) g(y) e^{-t(U + V)} = (\int f\,e^{-tU})(\int g\,
  e^{-tV})$.\footnote{In
  \texttt{Laplace/TwoD/AddSeparable.lean} as
  \texttt{integral\_separable\_addSeparable}.}
\item The 1D exact moments are also already in the seabed:
  \texttt{Laplace.OneD.quartic\_moment\_even} and
  \texttt{Laplace.OneD.sextic\_moment\_even}. T1 used them; Z1 uses
  them again, this time directly rather than via the
  \texttt{gibbsExpectation} layer.
\end{itemize}

This Tide is therefore the cleanest possible mirror of V2: the same
three-theorem packaging with different building blocks for the exact
reformulation step.

\section{The thing formalised}

Three theorems in \texttt{Laplace.TwoD}:

\begin{enumerate}[leftmargin=1.5em]
\item \textbf{Exact reformulation.} For $t > 0$ and $j, k \in
  \mathbb N$,
  $$
    \int\!\int x^{2j} y^{2k}\,e^{-tL(x, y)}\,dx\,dy
    = C'(j, k)\cdot t^{-((2j+1)/4 + (2k+1)/6)}
  $$
  with
  $$
    C'(j, k) := \tfrac{1}{6}\cdot 24^{(2j+1)/4}\cdot 720^{(2k+1)/6}
      \cdot \Gamma((2j+1)/4)\cdot \Gamma((2k+1)/6).
  $$
\item \textbf{Rescaled \texttt{Tendsto}.} As $t \to \infty$,
  $$
    t^{(2j+1)/4 + (2k+1)/6}\cdot \int\!\int x^{2j} y^{2k}\,e^{-tL}
    \to C'(j, k).
  $$
\item \textbf{\texttt{IsEquivalent}.} At $\mathrm{atTop}$,
  $$
    \int\!\int x^{2j} y^{2k}\,e^{-tL}
    \;\sim\; C'(j, k)\cdot t^{-((2j+1)/4 + (2k+1)/6)}.
  $$
\end{enumerate}

Named constant: \texttt{quarticSexticNumeratorConst}. File:
sibling to V2's moment-asymptotic file, in the same
\texttt{Laplace/TwoD/} folder.\footnote{Specifically
\texttt{Laplace/TwoD/QuarticSexticNumeratorAsymptotic.lean}, alongside
\texttt{QuarticSexticMomentAsymptotic.lean} from V2.} Added to
project root.

\section{Proof strategy}

The exact reformulation pipeline:
\begin{enumerate}[leftmargin=1.5em]
\item Apply \texttt{integral\_separable\_addSeparable} with
  $f = x \mapsto x^{2j}$, $g = y \mapsto y^{2k}$, integrability
  witnesses \texttt{quartic\_integrable\_pow\_pot} and
  \texttt{sextic\_integrable\_pow\_pot}. This factors the 2D integral
  into a product of two 1D integrals, but with the potentials still
  in \texttt{quarticPotential} / \texttt{sexticPotential} form.
\item Rewrite each 1D integrand from \texttt{quarticPotential x} to
  the explicit \texttt{x\textasciicircum 4 / 24} form (and similarly
  for sextic), via the \texttt{\_apply} \texttt{@[simp]} lemmas.
  Each is a small \texttt{congr 1; ext; ...; ring} block; mirror of
  T1's template.
\item Apply \texttt{quartic\_moment\_even} and
  \texttt{sextic\_moment\_even} to replace each 1D integral with its
  $(\text{const} \times (24/t)^{\text{exp}} \times \Gamma)$ form.
\item Unfold \texttt{quarticSexticNumeratorConst}, then standard
  \texttt{Real.rpow} algebra: \texttt{Real.div\_rpow} splits
  $(24/t)^a$ into $24^a / t^a$, \texttt{Real.rpow\_add} (valid for
  $t > 0$) combines $t^{-a} \cdot t^{-b} = t^{-(a + b)}$, and
  \texttt{Real.rpow\_neg} bridges negative exponents.
  \texttt{ring} closes.
\end{enumerate}

The two asymptotic forms (\texttt{Tendsto} and
\texttt{IsEquivalent}) follow from the exact reformulation by the
same eventual-equality idiom V2 used. Both reduce to ``the rescaled
function is eventually equal to a constant'' on the positive reals.

\section{Roadblocks and resolutions}

No roadblocks. The proof compiled green on first try. The lessons
from V2 (no \texttt{field\_simp} after explicit \texttt{rpow\_add}
normalisation; use \texttt{ring} at the bottom; \texttt{tendsto\_congr'}
+ eventual equality for the asymptotic packagings) carried directly.

The structural insight made before writing any code --- that the
seabed already provides the 2D-numerator factor lemma --- saved
roughly a third of the expected work (the original Z1 estimate
budgeted for proving a numerator factor lemma alongside,
${\sim}30$ lines).

\section{What was Mathlib, what was new}

\paragraph{From Mathlib.} The standard \texttt{Real.rpow} algebra
suite,\footnote{\texttt{Real.div\_rpow}, \texttt{Real.rpow\_add},
\texttt{Real.rpow\_neg}, plus \texttt{ring}.} the
\texttt{Tendsto}-to-constant idiom,\footnote{\texttt{tendsto\_const\_nhds}
and \texttt{Tendsto.congr'}.} and the
\texttt{Asymptotics.IsEquivalent} refl + congruence
pair.\footnote{\texttt{Asymptotics.IsEquivalent.refl} and
\texttt{IsEquivalent.congr\_left}.} Same Mathlib API V2 used.

\paragraph{From the seabed.} Five seabed lemmas, all proved on prior
tides: the 2D-numerator factor theorem from A2's
\texttt{AddSeparable.lean}, two 1D integrability witnesses, two 1D
\texttt{\_moment\_even} closed forms, and the two
\texttt{\_Potential\_apply} \texttt{@[simp]}
unfolds.\footnote{Specifically
\texttt{integral\_separable\_addSeparable},
\texttt{quartic\_integrable\_pow\_pot},
\texttt{sextic\_integrable\_pow\_pot},
\texttt{quartic\_moment\_even}, \texttt{sextic\_moment\_even},
\texttt{quarticPotential\_apply}, and
\texttt{sexticPotential\_apply}.} Nothing else from the laplace
seabed was load-bearing.

\paragraph{New here.} Three theorems plus the named constant
\texttt{quarticSexticNumeratorConst}. Architecturally identical to
V2; the value is in completing the numerator-vs-moment parity.

\section{Lessons}

\paragraph{Mirror-shaped tides are first-class.} A tide whose
candidate is structurally identical to a just-landed predecessor
\emph{should} skip the GPT consult --- the deliberation has already
happened, and re-running it adds token cost without changing the
output. The tide skill's preflight already names this path
(``proceed-without-GPT''); Z1 is a clean exemplar. The discipline
that justifies the skip is (a) a numerical sanity check covering at
least 5--10 test points, and (b) explicit acknowledgement in the
tide log that the candidate is mirror-shaped against the named
predecessor.

\paragraph{Seabed dig before deliberation pays off.} A 5-minute
\texttt{grep} of \texttt{AddSeparable.lean} surfaced the existing
2D-numerator factor lemma, which the survey direction had guessed
needed to be proved alongside Z1. Knowing this before drafting
candidates cut the Tide cost estimate by a third and made the
``mirror of V2'' framing exact.

\paragraph{Moment vs numerator scaling differ by the partition.}
V2's moment scales as $t^{-(j/2 + k/3)}$; Z1's numerator scales as
$t^{-((2j+1)/4 + (2k+1)/6)}$. The difference is exactly
$t^{-5/12} = t^{-1/4 - 1/6}$, which is the partition-function
scaling (quartic partition $\sim t^{-1/4}$, sextic
$\sim t^{-1/6}$). Worth recording as a sanity-check identity for
future asymptotic packagings.

\section{Follow-ups}

\begin{itemize}[leftmargin=1.5em]
\item \emph{2D partition function asymptotic.} The natural completing
  third object: the 2D partition function $\int\!\int e^{-tL}$ scales
  as $C''\cdot t^{-1/4 - 1/6}$. Same three-theorem packaging, even
  smaller proof (no polynomial weights, just the partition product).
  ${\sim}50$ lines.
\item \emph{Odd-numerator analogues.} For any odd index, the 2D
  numerator with $x^{2j+1}$ or $y^{2k+1}$ weight is identically zero
  by V1. The corresponding ``asymptotic packaging'' is the trivial
  $0 \sim 0$; skip.
\item \emph{Generic $k$-th-power potential 1D and 2D analogues.}
  The same template (1D moment closed form + 2D separability) works
  for $x^{2k}/(2k)!$ pure-monomial potentials in general. The 1D
  generic-$k$ asymptotic is exactly N1; extending to 2D follows the
  Z1 template. Out of scope for this tide.
\item \emph{Promote \texttt{quarticSexticNumeratorConst} to
  \texttt{lean-common}.} Only with a third consumer; currently the
  named constant has one use site.
\end{itemize}

\section*{Acknowledgements}

No GPT consult was used for Z1 (mirror-shaped tide off V2 per the
tide skill's preflight). The deliberation work was inherited from
V2's GPT consult, which set the three-theorem packaging shape and
the \texttt{tendsto\_congr'} idiom; both apply unchanged here. The
numerical sanity check at ten $(j, k, t)$ values confirmed the
$C'(j, k)$ formula to relative error ${\sim}10^{-16}$ before the
proof was written, and is preserved in the primer's tide-log
GPT-responses folder alongside this tide's markdown log.

\end{document}
